% Theory-refining H6 addendum generated from the capability-conversion redesign.
\subsubsection{Capability-Conversion Threshold}

The aggregate resilience effect of productive complexity may be non-linear. A country can develop isolated sophisticated export activities before it has the broad supplier, labor, and adjacent-product systems needed to convert those activities into economy-wide adjustment capacity. In this intermediate state, greater complexity may increase exposure to demanding international production networks without yet providing sufficient redundancy, substitutability, or reallocation capacity. Once a broader capability base is present, additional complexity should become more convertible into export recovery.

\begin{quote}
H6: Following tariff-induced disruption, increases in productive complexity are associated with weaker export recovery among countries below a minimum level of baseline productive complexity; this negative relationship attenuates once countries possess a sufficiently broad productive capability base.
\end{quote}

H6 is theory-refining rather than confirmatory because the motivating low-ECI pattern was identified in the current dataset. We operationalize baseline productive complexity as the 2015--2017 average ECI and estimate a pooled segmented country and year fixed-effects model with continuous pre-shock US--China nexus exposure. The selected breakpoint is -0.901 (wild-bootstrap grid interval -0.901 to 0.752). The primary outcome is log export recovery, with ratio-based, winsorized, pretrend-adjusted, small-economy exclusion, export-size-weighted, and alternative-exposure specifications reported as robustness checks.

The decisive tests are a negative low-regime slope and a positive difference between the high- and low-regime slopes. We report 999-replication country wild-bootstrap p-values, Benjamini--Hochberg adjusted q-values over the complete reported H6 exploration family, and repeated stratified country-holdout validation.